% SEGMENT OLP-0495-S01
% part: intuitionistic-logic
% chapter: introduction
% section: bhk-interpretation

\documentclass[../../../include/open-logic-section]{subfiles}

\begin{document}

\olfileid{int}{int}{bhk}

\olsection{பிரௌவர்--ஹெய்ட்டிங்--கொல்மொகோரோவ் பொருள்கோள்}

\begin{editorial}
  BHK பொருள்கோளைப் பயன்படுத்தி உள்ளுணர்வுவாதக் கூற்றுகளின்
  செல்லுபடிநிலையை நிறுவுவது குழப்பமானது; இதனை இன்னும் தெளிவாக விளக்க
  வேண்டும்.
\end{editorial}

உள்ளுணர்வுவாத இணைப்பிகளுக்கு முறைசாரா கட்டுமானப் பொருள்கோள் ஒன்று உள்ளது;
அது பொதுவாக பிரௌவர்--ஹெய்ட்டிங்--கொல்மொகோரோவ் பொருள்கோள் எனப்படுகிறது.
``ஆக்கம்'' என்ற கருத்துருவை அது பயன்படுத்துகிறது; இதைக் கட்டுமான
நிரூபிப்பாகக் கருதலாம். (முறைசாரான !!{derivation} அமைப்பிலுள்ள
!!a{derivation} என்ற கருத்துருவுடன் குழம்பாதிருக்க, BHK பொருள்கோளில்
``நிரூபிப்பு'' என்ற சொல்லைப் பயன்படுத்துவதில்லை.) இந்த உள்ளுணர்வுசார்
கருத்துருவை அடிப்படையாகக் கொண்டு, உள்ளுணர்வுவாத இணைப்பிகளின் பொருள்களை BHK
பொருள்கோள் விளக்குகிறது.

\begin{enumerate}
\item ஓர் அணுக் கூற்றின் ஆக்கம் எது என்பதை நாம் அறிவதாகக் கொள்கிறோம்.
\item $!A_1 \land !A_2$ இன் ஆக்கம் என்பது $\tuple{M_1, M_2}$ என்ற சோடி;
  இதில் $M_1$ ஆனது~$!A_1$ இன் ஆக்கம், $M_2$ ஆனது~$!A_2$ இன் ஆக்கம்.
\item $!A_1 \lor !A_2$ இன் ஆக்கம் என்பது $\tuple{s, M}$ என்ற சோடி;
  இதில் $s$ ஆனது~$1$ ஆகவும் $M$ ஆனது~$!A_1$ இன் ஆக்கமாகவும் இருக்கும்,
  அல்லது $s$ ஆனது~$2$ ஆகவும் $M$ ஆனது~$!A_2$ இன் ஆக்கமாகவும் இருக்கும்.
\item $!A \lif !B$ இன் ஆக்கம் என்பது~$!A$ இன் ஆக்கத்தை~$!B$ இன்
  ஆக்கமாக மாற்றும் ஒரு சார்பு.
\item $\lfalse$ (பொருத்தமின்மை) க்கு ஆக்கம் இல்லை.
\item $\lnot !A$ என்பது $!A \lif \lfalse$ என்பதன் ஒருபொருட்சொல்லாக
  வரையறுக்கப்படுகிறது. அதாவது, $\lnot !A$ இன் ஆக்கம் என்பது~$!A$ இன்
  ஆக்கத்தை~$\lfalse$ இன் ஆக்கமாக மாற்றும் ஒரு சார்பு.
\end{enumerate}

\begin{ex}
எடுத்துக்காட்டாக $\lnot \lfalse$ ஐ எடுத்துக்கொள்க. $\lfalse$ இன் எந்த
ஆக்கத்தையும் உள்ளீடாகப் பெற்று, $\lfalse$ இன் ஓர் ஆக்கத்தை வெளியீடாகத் தரும்
ஒரு சார்பே அதன் ஆக்கம். தெளிவாக, முற்றொருமைச் சார்பு~$\Id{}$ அத்தகைய ஓர்
ஆக்கம்: $\lfalse$ இன் ஆக்கம்~$M$ கொடுக்கப்பட்டால், $\Id{}(M) = M$ ஆனது
$\lfalse$ இன் ஓர் ஆக்கத்தைத் தருகிறது.
\end{ex}

பொதுவாகச் சொன்னால், $\lnot !A$ என்பது ``$!A$ இன் ஆக்கம் இயலாதது'' என்று
பொருள்படும்.

\begin{ex}
எந்தக் கூற்று~$!A$ க்கும் $!A \lif \lnot \lnot !A$ ஐ நிறுவுவோம்; இது
$!A \lif ((!A \lif \lfalse) \lif \lfalse)$ ஆகும். அதன் ஆக்கம் ஒரு
சார்பு~$f$ ஆக இருக்க வேண்டும்; $!A$ இன் ஆக்கம்~$M$ கொடுக்கப்பட்டால், அது
$(!A \lif \lfalse) \lif \lfalse$ இன் ஆக்கம்~$f(M)$ ஐத் திருப்பித் தர
வேண்டும். $(!A \lif \lfalse) \lif \lfalse$ இன் ஆக்கத்தை~$f$ எவ்வாறு
அமைக்கிறது என்பது இதோ: $!A \lif \lfalse$ இன் ஆக்கம்~$h$ உள்ளீடாகக்
கொடுக்கப்படும்போது, $\lfalse$ இன் ஓர் ஆக்கத்தை வெளியிடும் சார்பு $g$ ஒன்றை
வரையறுக்க வேண்டும். $g$ ஐப் பின்வருமாறு வரையறுக்கலாம்: உள்ளீடு~$h$ ஐ,
முன்னரே நாம் பெற்ற~$!A$ இன் ஆக்கம்~$M$ க்குப் பயன்படுத்துக. $h$ இன்
வெளியீடு~$h(M)$ ஆனது $\lfalse$ இன் ஆக்கம் என்பதால், $M$ ஆனது~$!A$ இன்
ஆக்கமாக இருந்தால் $f(M)(h) = h(M)$ ஆனது~$\lfalse$ இன் ஓர் ஆக்கமாகும்.
\end{ex}

\begin{ex}
$\lnot(!A \land \lnot !A)$ க்கு, அதாவது $(!A \land (!A \lif
\lfalse)) \lif \lfalse$ க்கு ஓர் ஆக்கத்தைத் தருவோம். இது சார்பு~$f$;
$!A \land (!A \lif \lfalse)$ இன் ஆக்கம்~$M$ உள்ளீடாகக் கொடுக்கப்பட்டால்,
அது~$\lfalse$ இன் ஓர் ஆக்கத்தைத் தருகிறது. இணைப்பு $!B_1 \land !B_2$ இன்
ஆக்கம் என்பது $\tuple{N_1, N_2}$ என்ற சோடி; இதில் $N_1$ ஆனது~$!B_1$ இன்
ஆக்கம், $N_2$ ஆனது~$!B_2$ இன் ஆக்கம். $!B_1 \land !B_2$ இன் ஓர்
ஆக்கத்திலிருந்து முறையே $!B_1$, $!B_2$ ஆகியவற்றின் ஆக்கங்களை மீட்டெடுக்கும்
சார்புகள் $p_1$, $p_2$ ஆகியவற்றைப் பின்வருமாறு வரையறுக்கலாம்:
\begin{align*}
  p_1(\tuple{N_1, N_2}) &= N_1\\
  p_2(\tuple{N_1, N_2}) &= N_2
\end{align*}
$f$ செய்வது இதுதான்: முதலில் அது தனது உள்ளீடு~$M$ க்கு $p_1$ ஐப்
பயன்படுத்துகிறது. இதனால்~$!A$ இன் ஓர் ஆக்கம் கிடைக்கிறது. பின்னர் $p_2$ ஐ
$M$ க்குப் பயன்படுத்துவதால் $!A \lif \lfalse$ இன் ஓர் ஆக்கம் கிடைக்கிறது.
அத்தகைய ஆக்கம், $!A$ இன் ஓர் ஆக்கத்தை உள்ளீடாகக் கொடுத்தால் $\lfalse$ இன்
ஓர் ஆக்கத்தைத் தரும் சார்பு~$p_2(M)$ ஆகும். வேறு சொற்களில், $p_2(M)$ ஐ
$p_1(M)$ க்குப் பயன்படுத்தினால், $\lfalse$ இன் ஓர் ஆக்கம் கிடைக்கிறது.
எனவே $f(M) = p_2(M)(p_1(M))$ என வரையறுக்கலாம்.
\end{ex}

\begin{ex}
$((!A \land !B) \lif !C) \lif (!A \lif (!B \lif !C))$ இன் ஓர்
ஆக்கத்தைத் தருவோம்; அதாவது $(!A \land !B) \lif !C$ இன் ஆக்கம்~$g$ ஐ
$(!A \lif (!B \lif !C))$ இன் ஓர் ஆக்கமாக மாற்றும் சார்பு~$f$ ஐத் தருவோம்.
ஆக்கம்~$g$ தானே ஒரு சார்பு ($!A \land !B$ இன் ஆக்கங்களிலிருந்து~$!C$ இன்
ஆக்கங்களுக்குச் செல்லும் சார்பு). மேலும் வெளியீடு~$f(g)$ என்பது~$!A$ இன்
ஆக்கங்களிலிருந்து,~$!B$ இன் ஆக்கங்களிலிருந்து~$!C$ இன் ஆக்கங்களுக்குச்
செல்லும் சார்புகளுக்குச் செல்லும் சார்பு~$h_g$ ஆகும்.
% SOURCE CORRECTION TA-IL-001: restore the formula metavariable marker on C.

சரி, இது குழப்பமாக இருக்கிறது. உள்ளீடு~$g$ க்கான $f$ இன் வெளியீடாக இருக்கும்
ஒரு குறிப்பிட்ட சார்பு~$h_g$ ஐ அமைக்க வேண்டும். $h_g$ இன் உள்ளீடு~$!A$ இன்
ஆக்கம்~$M$. $h_g(M)$ இன் வெளியீடு,~$!B$ இன் ஆக்கங்கள்~$N$ இலிருந்து~$!C$ இன்
ஆக்கங்களுக்குச் செல்லும் சார்பு~$k_M$ ஆக இருக்க வேண்டும். $k_{g,M}(N) =
g(\tuple{M, N})$ எனக் கொள்க. $\tuple{M, N}$ ஆனது~$!A \land !B$ இன் ஓர்
ஆக்கம் என்பதை நினைவுகூர்க. எனவே $k_{g,M}$ ஆனது $!B \lif !C$ இன் ஓர்
ஆக்கம்:~$!B$ இன் ஆக்கங்கள்~$N$ ஐ~$!C$ இன் ஆக்கங்களாக அது மாற்றுகிறது.
இப்போது $h_g(M) = k_{g,M}$ எனக் கொள்க. இது~$!A$ இன் ஆக்கங்கள்~$M$ ஐ
$!B \lif !C$ இன் ஆக்கங்கள்~$k_{g,M}$ ஆக மாற்றும் சார்பு. இப்போது $f(g) =
h_g$ எனக் கொள்க. இது $(!A \land !B) \lif !C$ இன் ஆக்கங்கள்~$g$ ஐ
$!A \lif (!B \lif !C)$ இன் ஆக்கங்களாக மாற்றும் சார்பு. அப்பாடா!{}
\end{ex}

$!A \lor \lnot !A$ என்ற கூற்று விலக்கப்பட்ட நடுவின் விதி எனப்படுகிறது.
குறிப்பிட்ட சில~$!A$ க்களுக்கு இதை நிறுவலாம் (எடுத்துக்காட்டாக, $\lfalse
\lor \lnot \lfalse$); ஆனால் பொதுவாக நிறுவ முடியாது. ஏனெனில்
உள்ளுணர்வுவாதப் பிரிப்புக்கு அதன் பிரிப்புறுப்புகளில் ஒன்றின் ஆக்கம் தேவை;
ஆனால் தற்போது நிறுவவும் மறுக்கவும் இயலாத கூற்றுகள் உள்ளன (கோல்ட்பாக்
கருதுகோள் போன்றவை). எனினும், விலக்கப்பட்ட நடுவின் விதியை மறுக்கவும்
முடியாது: அதாவது, $\lnot \lnot (!A \lor \lnot !A)$ மெய்யாகிறது.

\begin{ex}
$\lnot \lnot (!A \lor \lnot !A)$ ஐ நிறுவ, $\lnot(!A \lor \lnot !A)$
இன், அதாவது $(!A \lor (!A \lif \lfalse)) \lif \lfalse$ இன் ஓர்
ஆக்கத்தை~$\lfalse$ இன் ஓர் ஆக்கமாக மாற்றும் சார்பு~$f$ நமக்குத் தேவை.
வேறு சொற்களில், $g$ ஆனது $\lnot(!A \lor \lnot !A)$ இன் ஆக்கமாக இருந்தால்,
$f(g)$ ஆனது~$\lfalse$ இன் ஆக்கமாக இருக்குமாறு ஒரு சார்பு~$f$ தேவை.

$g$ ஆனது $\lnot(!A \lor \lnot !A)$ இன் ஓர் ஆக்கம் என்று கொள்க; அதாவது,
$!A \lor \lnot !A$ இன் ஓர் ஆக்கத்தை~$\lfalse$ இன் ஓர் ஆக்கமாக மாற்றும்
சார்பு அது. $!A \lor \lnot !A$ இன் ஆக்கம் என்பது $\tuple{s, M}$ என்ற சோடி;
இதில் $s = 1$ ஆகவும் $M$ ஆனது~$!A$ இன் ஆக்கமாகவும் இருக்கும், அல்லது $s = 2$
ஆகவும் $M$ ஆனது~$\lnot !A$ இன் ஆக்கமாகவும் இருக்கும். $h_1$ என்பது~$!A$ இன்
ஆக்கம்~$M_1$ ஐ $!A \lor \lnot !A$ இன் ஓர் ஆக்கமாக மாற்றும் சார்பாக
இருக்கட்டும்: அது $M_1$ ஐ $\tuple{1, M_1}$ க்கு அனுப்புகிறது. மேலும் $h_2$
என்பது~$\lnot !A$ இன் ஆக்கம்~$M_2$ ஐ $!A \lor \lnot !A$ இன் ஓர் ஆக்கமாக
மாற்றும் சார்பாக இருக்கட்டும்: அது $M_2$ ஐ $\tuple{2, M_2}$ க்கு அனுப்புகிறது.
% SOURCE CORRECTION TA-IL-002: h_1 must inject its own input M_1.

$k$ ஆனது $\comp{h_1}{g}$ ஆக இருக்கட்டும்:~$!A$ இன் ஓர் ஆக்கம்
கொடுக்கப்பட்டால், $\lfalse$ இன் ஓர் ஆக்கத்தைத் திருப்பித் தரும் சார்பு இது;
அதாவது $!A \lif \lfalse$ அல்லது $\lnot !A$ இன் ஓர் ஆக்கம். இப்போது $l$
ஆனது $\comp{h_2}{g}$ ஆக இருக்கட்டும். $\lnot !A$ இன் ஓர் ஆக்கம்
கொடுக்கப்பட்டால், $\lfalse$ இன் ஓர் ஆக்கத்தைத் தரும் சார்பு இது. $k$ ஆனது
$\lnot !A$ இன் ஆக்கம் என்பதால், $l(k)$ ஆனது~$\lfalse$ இன் ஓர் ஆக்கம்.

மொத்தத்தில், $\lnot(!A \lor \lnot !A)$ இன் ஆக்கம்~$g$ ஐ~$\lfalse$ இன் ஓர்
ஆக்கமாக எவ்வாறு மாற்றலாம் என்பதை விவரித்துள்ளோம்; அதாவது $\lnot(!A \lor
\lnot !A)$ இன் ஆக்கம்~$g$ ஐ~$\lfalse$ இன் ஆக்கம்~$l(k)$ க்கு அனுப்பும்
சார்பு~$f$, $\lnot\lnot(!A \lor \lnot !A)$ இன் ஓர் ஆக்கமாகும்.
\end{ex}

நீங்கள் காண்பதுபோல், !!{formula}s இன் உள்ளுணர்வுவாதச் செல்லுபடிநிலையைக்
காட்ட BHK பொருள்கோளைப் பயன்படுத்துவது விரைவிலேயே சுமையானதாகவும்
குழப்பமானதாகவும் ஆகிறது. நல்வாய்ப்பாக, உள்ளுணர்வுவாதத் தருக்கத்திற்கு மேலும்
சிறந்த !!{derivation} அமைப்புகளும் மேலும் துல்லியமான பொருண்மையியல்
பொருள்கோள்களும் உள்ளன.
\end{document}
